\section{Composition series}\label{sec:composition_series}

\paragraph{Normal series}

\begin{definition}\label{def:normal_series}\mcite[95]{Курош1967ТеорияГрупп}
  A \term[bg=нормален ред (\cite[def. XII.6]{ГеновМиховскиМоллов1991Алгебра}), ru=нормальный ряд, en=normal series (\cite[207]{Aluffi2009Algebra})]{normal series} of \hyperref[def:group/subobject]{subgroups} of \( G \) is a strictly decreasing finite sequence
  \begin{equation}\label{eq:def:normal_series}
    \set{ e } = G_n \nsubseteq G_{n-1} \nsubseteq \cdots \nsubseteq G_1 \nsubseteq G_0 = G,
  \end{equation}
  where, for each index \( k \), \( G_{k+1} \) is normal in \( G_k \) (but not necessarily in \( G \)).

  We associate with each normal series the sequence of quotients \( G_{k+1} / G_k \).

  \begin{thmenum}
    \thmitem{def:normal_series/refinement} If one normal series for \( G \) is a \hyperref[def:subsequence]{subsequence} of another, we call the second one a \term[ru=уплотнение, en=refinement (\cite[197]{Rotman2015AdvancedModernAlgebraPart1})]{refinement} of the first.

    \thmitem{def:normal_series/equivalence}\mcite[96]{Курош1967ТеорияГрупп} If two normal series for \( G \) have the same length and if their quotients are pairwise isomorphic, we say that the series are \term[ru=изоморфные (ряды), en=equivalent (series) (\cite[197]{Rotman2015AdvancedModernAlgebraPart1})]{equivalent}.

    \thmitem{def:normal_series/composition}\mcite[97]{Курош1967ТеорияГрупп} The \hyperref[def:extremal_points/maximal_and_minimal_element]{maximal} normal series for \( G \) with respect to refinement are called \term[ru=композиционный ряд, en=composition series (\cite[302]{Rotman2015AdvancedModernAlgebraPart1})]{composition series}.
  \end{thmenum}
\end{definition}
\begin{comments}
  \item There are some inconsistencies regarding this definition.
  \begin{itemize}
    \item We follow the presentation of \textcite[96]{Курош1967ТеорияГрупп}.
    \item \Textcite[205]{Aluffi2009Algebra} allows a normal series to be infinitely descending.
    \item \Textcite[244]{Jacobson1985BasicAlgebraI}, \textcite[192]{Rotman2015AdvancedModernAlgebraPart1} and \textcite[173]{Knapp2016BasicAlgebra} allow nonstrict inequalities in \eqref{eq:def:normal_series} for normal series, but not for composition series.

    Rotman defines the length of a normal series as the number of \emph{nontrivial} subgroups.

    \item In a footnote, \textcite[193]{Rotman2015AdvancedModernAlgebraPart1} hints that normal series can be called \enquote{\textru{субнормальная}} (\enquote{subnormal}) because \( G_k \) is only \hyperref[def:subnormal_group]{subnormal} (but not necessarily normal) in \( G \).
  \end{itemize}
\end{comments}

\begin{definition}\label{def:subnormal_group}\mcite[98]{Курош1967ТеорияГрупп}
  We say that a \hyperref[def:group/subobject]{subgroup} \( H \) of \( G \) is \term[ru=достижимая (подгруппа); субнормальная (подгруппа) (\cite[439]{Курош1967ТеорияГрупп}), en=subnormal (subgroup) (\cite[192]{Rotman2015AdvancedModernAlgebraPart1})]{subnormal} if \( H \) belongs to a \hyperref[def:normal_series]{normal series} for \( G \).
\end{definition}

\begin{example}\label{ex:def:normal_series}
  We list examples of \hyperref[def:normal_series]{normal series}:
  \begin{thmenum}
    \thmitem{ex:def:normal_series/integers} In the additive group \( (\BbbZ, +) \) of \hyperref[def:integers]{integers}, every integer \( n \) induces the normal series
    \begin{equation*}
      \set{ 0 } \subsetneq \braket{ n } \subsetneq \BbbZ.
    \end{equation*}

    We can refine this series via multiples of \( n \):
    \begin{equation*}
      \set{ 0 } \subsetneq \cdots \subsetneq \braket{ 3n } \subsetneq  \braket{ 2n } \subsetneq \braket{ n } \subsetneq \BbbZ.
    \end{equation*}

    We can keep refining the series, however we can never reach \( \set{ 0 } \) in this way. In fact, all nontrivial subgroups of \( \BbbZ \) are countably infinite, so we cannot build a composition series.

    More generally, every \hyperref[def:partial_order_chain]{chain} in the natural number divisibility lattice from \cref{thm:natural_number_divisibility_lattice} induces an \enquote{infinite normal series}.

    \thmitem{ex:def:normal_series/subnormal} We will study normal series in the \hyperref[def:dihedral_group]{Dihedral group}
    \begin{equation*}
      D_4 = \braket{ a, b \given a^4 = b^2 = e \T{and} bab = a^3 }.
    \end{equation*}

    \Cref{ex:d2_normal_subgroups} implies that the groups \( \braket{ a^2 } \) and \( \braket{ a^2 b } \) are normal, hence
    \begin{equation}\label{eq:ex:def:normal_series/subnormal}
      \set{ e } \subsetneq \braket{ a^2 } \subsetneq \braket{ a^2, b } \subsetneq D_4
    \end{equation}
    is a normal series. Since the quotients have index \( 2 \), this series is compositional.

    Another subgroup discussed in \cref{ex:d2_normal_subgroups} is \( \braket{ a^2 b } \), which is not normal in \( D_4 \). It is normal in \( \braket{ a^2, b } \), however, because it has index \( 2 \) there. Thus, \( \braket{ a^2 b } \) is an example of a \hyperref[def:subnormal_group]{subnormal group} that is not normal.

    This leads to another compositional series:
    \begin{equation}
      \set{ e } \subsetneq \braket{ a^2 b } \subsetneq \braket{ a^2, b } \subsetneq D_4.
    \end{equation}
  \end{thmenum}
\end{example}

\paragraph{Commutator subgroups}

\begin{definition}\label{def:group_commutator}\mcite[245]{Jacobson1985BasicAlgebraI}
  Let \( G \) be an arbitrary group. We define the \term[ru=коммутатор (\cite[37]{Кострикин2001АлгебраЧасть3})]{commutator} of the elements \( x \) and \( y \) as
  \begin{equation*}
    [x, y] \coloneqq \underbrace{xy(yx)^{-1}}_{xyx^{-1}y^{-1}}.
  \end{equation*}

  We call the subgroup \hyperref[def:group/subobject]{generated} by all the commutators in \( G \) the \term[ru=коммутант (группы) (\cite[37]{Кострикин2001АлгебраЧасть3})]{commutator subgroup} and denote it by \( [G, G] \).

  In analogy with \hyperref[def:differentiability]{function derivatives} and \hyperref[def:derived_set]{derived sets} in topology, we call \( [G, G] \) the \term[ru=производная подгруппа (\cite[37]{Кострикин2001АлгебраЧасть3})]{derived subgroup} and denote \( [G, G] \) by \( G' \). Via iteration, we extend this notation to \( G^{(k)} \).
\end{definition}
\begin{comments}
  \item Compare this definition to ring commutators defined in \cref{def:ring_commutator}.
\end{comments}

\begin{proposition}\label{thm:commutator_subgroup_is_normal}
  The \hyperref[def:group_commutator]{commutator subgroup} of any group is \hyperref[def:normal_subgroup]{normal}.
\end{proposition}
\begin{proof}
  Let \( G \) be an arbitrary group. Fix \( g \) from \( G \) and \( n \) from \( [G, G] \). Via \fullref{thm:induction_on_generated_substructures}, we will show that the conjugate \( g n g^{-1} \) belongs to \( [G, G] \).

  \begin{itemize}
    \item If \( n = [x, y] \), then
    \begin{equation*}
      g n g^{-1}
      =
      g(xyx^{-1}y^{-1})g^{-1}
      =
      (gxg^{-1}) (gyg^{-1}) (\underbrace{gx^{-1}g^{-1}}_{(gxg^{-1})^{-1}}) (\underbrace{gy^{-1}g^{-1}}_{(gyg^{-1})^{-1}})
      =
      [gxg^{-1}, gyg^{-1}].
    \end{equation*}

    Hence, \( gng^{-1} \) is a commutator and thus also belongs to \( [G, G] \).

    \item Suppose that \( n = n_1 \cdot n_2 \) for some elements \( n_1 \) and \( n_2 \) of \( [G, G] \) and suppose that the inductive hypothesis holds for \( n_1 \) and \( n_2 \).

    Thus, since \( n_1 \) and \( n_2 \) belong to \( N \), then \( g n_1 g^{-1} \) and \( g n_2 g^{-1} \) also do. Then so does
    \begin{equation*}
      gng^{-1}
      =
      g n_1 n_2 g^{-1}
      =
      g n_1 g^{-1} \cdot g n_2 g^{-1}.
    \end{equation*}
  \end{itemize}
\end{proof}

\begin{proposition}\label{thm:quotient_is_abelian_iff_subgroup_contains_commutator}
  The \hyperref[def:group/quotient]{quotient} \( G / N \) of a \hyperref[def:group]{group} is \hyperref[def:abelian_group]{abelian} if and only if \( N \) contains the \hyperref[def:group_commutator]{commutator subgroup} \( [G, G] \) of \( G \).
\end{proposition}
\begin{comments}
  \item Compare this to \cref{thm:quotient_is_abelian_iff_ideal_contains_commutator} for ring commutators.
\end{comments}
\begin{proof}
  \SufficiencySubProof Suppose that \( G / N \) is abelian. Fix two elements \( x \) and \( y \) of \( G \).

  Due to commutativity, we have
  \begin{equation*}
    xyN = xN \cdot yN = yN \cdot xN = yxN.
  \end{equation*}

  Then
  \begin{equation*}
    N = xyN \cdot (yxN)^{-1} = (xy(yx)^{-1})N,
  \end{equation*}
  which implies that the commutator \( xy(yx)^{-1} \) belongs to \( N \).

  Since \( x \) and \( y \) were arbitrary, we conclude that the commutator subgroup \( [G, G] \) belongs to \( N \).

  \NecessitySubProof Suppose that \( [G, G] \subseteq N \). Fix two elements \( x \) and \( y \) of \( G \).

  We have
  \begin{equation*}
    xy = \underbrace{xy (yx)^{-1}}_{[x, y]} yx,
  \end{equation*}
  hence \( xy \) belongs to
  \begin{equation*}
    Nyx
    \reloset {\ref{thm:normal_subgroup_left_right_cosets}} =
    yxN
    =
    yN \cdot xN.
  \end{equation*}

  But
  \begin{equation*}
    xy \in xyN = xN \cdot yN.
  \end{equation*}

  The cosets \( xyN \) and \( yxN \) intersect, and, as equivalence classes, they can only be equal.

  Therefore,
  \begin{equation*}
    xN \cdot yN = yN \cdot xN,
  \end{equation*}
  which shows that \( G / N \) is abelian.
\end{proof}

\begin{example}\label{ex:def:group_commutator}
  We list some examples of \hyperref[def:group_commutator]{commutators} in groups:
  \begin{thmenum}
    \thmitem{ex:def:group_commutator/gl2}\mcite{MathSE:commutator_group_of_gl2_is_sl2} Consider the \hyperref[def:linear_groups]{general linear group} \( \grp{GL}_2(\BbbK) \) over any \hyperref[def:field]{field} of \hyperref[def:ring_characteristic]{characteristic} different from \( 2 \).

    We will show that its commutator subgroup is the special linear group \( \grp{SL}_2(\BbbK) \). This is obvious in one direction --- for any pair of matrices \( A \) and \( B \), we have
    \begin{equation*}
      \det(AB A^{-1} B^{-1}) = \det(A) \det(B) \det(A)^{-1} \det(B)^{-1} = 1,
    \end{equation*}
    thus the commutators belongs to \( \grp{SL}_n(\BbbK) \).

    In the other direction, we will use \cref{thm:transversion_matrices_generate_sln}, which states that \( \grp{SL}_n(\BbbK) \) is generated by the family of transvection matrices (parameterized by \( x \))
    \begin{align*}
      \begin{pmatrix}
        1 & x \\
        0 & 1
      \end{pmatrix}
      &&
      \T{and}
      &&
      \begin{pmatrix}
        1 & 0 \\
        x & 1
      \end{pmatrix}.
    \end{align*}

    First, note that
    \begin{align*}
      \begin{pmatrix}
        1 & x \\
        0 & 1
      \end{pmatrix}
      =
      \bracks*
        {
          \begin{pmatrix}
            2 & 0 \\
            0 & 1
          \end{pmatrix},
          \begin{pmatrix}
            1 & x \\
            0 & 1
          \end{pmatrix}
        }.
    \end{align*}

    Indeed,
    \begin{equation*}
      \begin{pmatrix}
        2 & 0 \\
        0 & 1
      \end{pmatrix}
      \begin{pmatrix}
        1 & x \\
        0 & 1
      \end{pmatrix}
      =
      \begin{pmatrix}
        2 & 2x \\
        0 & 1
      \end{pmatrix}
    \end{equation*}
    and
    \begin{equation*}
      \begin{pmatrix}
        2 & 0 \\
        0 & 1
      \end{pmatrix}^{-1}
      \begin{pmatrix}
        1 & x \\
        0 & 1
      \end{pmatrix}^{-1}
      =
      \begin{pmatrix}
        1/2 & 0 \\
        0   & 1
      \end{pmatrix}
      \begin{pmatrix}
        1 & -x \\
        0 & 1
      \end{pmatrix}
      =
      \begin{pmatrix}
        1/2 & -x/2 \\
        0   & 1
      \end{pmatrix},
    \end{equation*}
    so
    \begin{equation*}
      \begin{pmatrix}
        2 & 0 \\
        0 & 1
      \end{pmatrix}
      \begin{pmatrix}
        1 & x \\
        0 & 1
      \end{pmatrix}
      \begin{pmatrix}
        2 & 0 \\
        0 & 1
      \end{pmatrix}^{-1}
      \begin{pmatrix}
        1 & x \\
        0 & 1
      \end{pmatrix}^{-1}
      =
      \begin{pmatrix}
        2 & 2x \\
        0 & 1
      \end{pmatrix}
      \begin{pmatrix}
        1/2 & -x/2 \\
        0   & 1
      \end{pmatrix}
      =
      \begin{pmatrix}
        1 & x \\
        0 & 1
      \end{pmatrix}.
    \end{equation*}

    By transposing the above, we can conclude that
    \begin{align*}
      \begin{pmatrix}
        1 & 0 \\
        x & 1
      \end{pmatrix}
      =
      \begin{pmatrix}
        1 & x \\
        0 & 1
      \end{pmatrix}^T
      =
      \bracks*
        {
          \begin{pmatrix}
            1 & 0 \\
            -x & 1
          \end{pmatrix},
          \begin{pmatrix}
            1/2 & 0 \\
            0   & 1
          \end{pmatrix}
        }.
    \end{align*}

    Therefore, both generators of \( \grp{SL}_2(\BbbK) \) are commutators, which makes \( \grp{SL}_2(\BbbK) \) a subgroup of the commutator subgroup of \( \grp{GL}_2(\BbbK) \).

    But we have already shown the converse inclusion. Therefore, the commutator subgroup of \( \grp{GL}_2(\BbbK) \) coincides with \( \grp{SL}_2(\BbbK) \).

    \thmitem{ex:def:group_commutator/s3} Consider the \hyperref[def:symmetric_group]{symmetric group} \( S_3 \) discussed in \cref{ex:s3}.

    To obtain its commutator subgroup, we must compute \( \sigma \rho \sigma^{-1} \rho^{-1} \) for any pair of permutations \( \sigma \) and \( \rho \), however this computation can be vastly simplified.

    First, note that transpositions are involutions, while the two \( 3 \)-cycles in \( S_3 \) are each other's inverses. In other words, the permutation \( \sigma \) is odd if and only if \( \sigma^{-1} \) is odd.

    \Cref{thm:permutation_product_parity} implies that \( \sigma \rho \) is odd if and only if \( \sigma \) and \( \rho \) have different signs odd, and similarly for \( \sigma^{-1} \rho^{-1} \). Therefore, the signs of \( \sigma \rho \) and \( \sigma^{-1} \rho^{-1} \) match.

    Applying \cref{thm:permutation_product_parity} again, we conclude that product \( (\sigma \rho) (\sigma^{-1} \rho^{-1}) \) is always an even permutation.

    Hence, every commutator in \( S_3 \) is an even permutation. We can easily show that every even permutation in \( S_3 \) occurs as a commutator:
    \begin{itemize}
      \item The identity \( \id \) occurs as the commutator of \( \cycle{ 1, 2 } \) and \( \cycle{ 1, 2 } \).
      \item \( \cycle{ 1, 2, 3 } \) occurs as the commutator of \( \cycle{ 1, 2 } \) and \( \cycle{ 1, 3 } \):
      \begin{equation*}
        \parens[\big]{ \cycle{ 1, 2 } \bincirc \cycle{ 1, 3 } } \bincirc \parens[\big]{ \cycle{ 1, 2 } \bincirc \cycle{ 1, 3 } }
        =
        \cycle{ 1, 3, 2 } \bincirc \cycle{ 1, 3, 2 }
        =
        \cycle{ 1, 2, 3 }.
      \end{equation*}

      \item \( \cycle{ 1, 3, 2 } \) occurs as the commutator of \( \cycle{ 1, 2 } \) and \( \cycle{ 2, 3 } \):
      \begin{equation*}
        \parens[\big]{ \cycle{ 1, 2 } \bincirc \cycle{ 2, 3 } } \bincirc \parens[\big]{ \cycle{ 1, 2 } \bincirc \cycle{ 2, 3 } }
        =
        \cycle{ 1, 2, 3 } \bincirc \cycle{ 1, 2, 3 }
        =
        \cycle{ 1, 3, 2 }.
      \end{equation*}
    \end{itemize}

    Therefore, the commutator subgroup of \( S_3 \) is simply the alternative group \( A_3 \).

    \thmitem{ex:def:group_commutator/dihedral} Consider the \hyperref[def:dihedral_group]{Dihedral group}
    \begin{equation*}
      D_n = \braket{ a, b \given a^n = b^2 = e \T{and} bab = a^{-1} }.
    \end{equation*}

    We will show that the commutator subgroup is \( \braket{ a^2 } \).

    First, we will prove that the commutator is a subgroup of \( \braket{ a^2 } \), i.e. that for any \( x \) and \( y \) in \( D_n \), the commutator \( xy x^{-1} y^{-1} \) is a power of \( a^2 \).

    \begin{itemize}
      \item If \( x = e \) or \( y = e \), then naturally \( xy x^{-1} y^{-1} = e \)
      \item If \( x = a^m \) and \( y = a^k \), then \( xy x^{-1} y^{-1} = a^{m+k-m-k} = a^0 = e \).
      \item If \( x = a^m b \) and \( y = a^k \), then
      \begin{align*}
        xyx^{-1} y^{-1}
        &=
        a^m b \cdot a^k \cdot b a^{-m} \cdot a^{-k}
        = \\ &=
        (a^m b a^k) b a^{-m-k}
        \reloset {\eqref{eq:thm:dihedral_group_identities/anbam}} = \\ &=
        a^{m-k} b b a^{-m-k}
        = \\ &=
        a^{m-k-m-k}
        = \\ &=
        a^{-2k}.
      \end{align*}

      \item Finally, if \( x = a^m b \) and \( y = a^k b \), then
      \begin{align*}
        xyx^{-1} y^{-1}
        &=
        a^m b \cdot a^k b \cdot b a^{-m} \cdot b a^{-k}
        = \\ &=
        a^m b a^{k-m} b a^{-k}
        \reloset {\eqref{eq:thm:dihedral_group_identities/anbam}} = \\ &=
        b a^{2k-2m} b
        \reloset {\eqref{eq:thm:dihedral_group_identities/ba2nb}} = \\ &=
        a^{-2m-2k}.
      \end{align*}
    \end{itemize}

    Incidentally, we have shown that \( [b, a^{-k}] = a^{2k} \), so any power of \( a^2 \) can be attained. This shows that the commutator of \( D_n \) coincides with with \( \braket{ a^2 } \).
  \end{thmenum}
\end{example}

\paragraph{Group abelianization}

\begin{proposition}\label{thm:abelian_group_iff_trivial_commutator}
  A group is \hyperref[def:abelian_group]{abelian} if and only if its \hyperref[def:group_commutator]{commutator} is trivial.
\end{proposition}
\begin{comments}
  \item More generally, we may take the quotient by the commutator to obtain the \hyperref[def:group_abelianization]{abelianization} of a group.
\end{comments}
\begin{proof}
  \SufficiencySubProof If \( G \) is abelian, then
  \begin{equation*}
    [x, y] = xyx^{-1}y^{-1} = (xx^{-1}) (yy^{-1}) = e.
  \end{equation*}

  \NecessitySubProof If \( [G, G] \) is trivial, then \( [x, y] = xyx^{-1}y^{-1} = e \), so \( xy = yx \).
\end{proof}

\begin{definition}\label{def:group_abelianization}\mcite[example 4.5.14(b)]{Leinster2014BasicCategoryTheory}
  We define the \term{abelianization} of a group \( G \) as its \hyperref[def:group/quotient]{quotient} \( G / [G, G] \) by its \hyperref[def:group_commutator]{commutator group} \( [G, G] \).
\end{definition}
\begin{comments}
  \item \Cref{thm:quotient_is_abelian_iff_subgroup_contains_commutator} implies that \( G / [G, G] \) is abelian, which justifies the name.
  \item Compare this ring group abelianization from \cref{def:ring_abelianization}.
\end{comments}
\begin{defproof}
  \Cref{thm:commutator_subgroup_is_normal} implies that \( [G, G] \) is normal, hence we are allowed to take quotients.
\end{defproof}

\begin{theorem}[Group abelianization universal property]\label{thm:group_abelianization_universal_property}
  The \hyperref[def:group_abelianization]{abelianization} \( G / [G, G] \) of a group \( G \) satisfies the following \hyperref[rem:universal_mapping_property]{universal mapping property}:
  \begin{displayquote}
    For every abelian group \( H \), every \hyperref[def:group/morphism]{group homomorphism} \( \varphi: G \to H \) \hyperref[def:factors_through]{uniquely factors through} \( G / [G, G] \). More precisely, there exists a unique group homomorphism \( \widetilde{\varphi}: G / [G, G] \to H \) such that the following diagram commutes:
    \begin{equation}\label{eq:thm:group_abelianization_universal_property/diagram}
      \begin{aligned}
        \includegraphics[page=1]{output/thm__group_abelianization_universal_property}
      \end{aligned}
    \end{equation}
  \end{displayquote}
\end{theorem}
\begin{comments}
  \item Via \cref{rem:universal_mapping_property}, the abelianization functor becomes \hyperref[def:category_adjunction]{left adjoint} to the \hyperref[def:concrete_category]{forgetful functor}
  \begin{equation*}
    U: \cat{Ab} \to \cat{Grp}.
  \end{equation*}

  \item Compare this result to \fullref{thm:ring_abelianization_universal_property} for ring abelianization.
\end{comments}
\begin{proof}
  Let \( H \) be an abelian group and let \( \varphi: G \to H \) be a group homomorphism.

  We want \( \oline{\varphi}: G / [G, G] \to H \) to satisfy
  \begin{equation}\label{eq:thm:group_abelianization_universal_property/homomorphism}
    \oline{\varphi}(\pi_G(x)) = \varphi(x).
  \end{equation}

  In order to use \eqref{eq:thm:group_abelianization_universal_property/homomorphism} as a definition, we must prove that \( x \cong y \pmod {[G, G]} \) implies that \( \varphi(x) = \varphi(y) \), that is, that \( xy^{-1} \in [G, G] \) implies \( \varphi(xy^{-1}) = e_H \). We will show via \fullref{thm:induction_on_generated_substructures} on \( n \in [G, G] \) that \( \varphi(n) = e_H \).
  \begin{itemize}
    \item If \( n = [x, y] \), then, by the commutativity of \( H \),
    \begin{equation*}
      \varphi(n)
      =
      \varphi(x) \varphi(y) \cdot \varphi(x)^{-1} \cdot \varphi(y)^{-1}
      =
      \parens[\big]{ \varphi(x) \cdot \varphi(x)^{-1} } \cdot \parens[\big]{ \varphi(y) \cdot \varphi(y)^{-1} }
      =
      e_H.
    \end{equation*}

    \item If \( n = n_1 \cdot n_2 \) for \( n_1 \) and \( n_2 \) from \( [G, G] \) and if the inductive hypothesis holds for \( n_1 \) and \( n_2 \), then
    \begin{equation*}
      \varphi(n)
      =
      \varphi(n_1) \cdot \varphi(n_2)
      =
      e_H \cdot e_H
      =
      e_H.
    \end{equation*}
  \end{itemize}

  Therefore, the function \( \oline{\varphi} \) is well-defined via \eqref{thm:group_abelianization_universal_property}.
\end{proof}

\paragraph{Solvable groups}

\begin{definition}\label{def:derived_group_series}\mcite[def. IV.3.9]{Aluffi2009Algebra}
  We define the \term[ru=производный ряд (\cite[48]{Кострикин2001АлгебраЧасть3})]{derived series} of the \hyperref[def:group]{group} \( G \) as the descending sequence of \hyperref[def:group_commutator]{commutator subgroups}
  \begin{equation*}
    \cdots \subseteq G^{(k)} \subseteq \cdots \subseteq G' \subseteq G.
  \end{equation*}
\end{definition}
\begin{comments}
  \item This is nor necessarily a normal series in the sense of \cref{def:normal_series} because it is not strictly descending and, furthermore, it may not stabilize (see \cref{thm:simple_solvable_group_is_abelian}).
\end{comments}

\begin{definition}\label{def:solvable_group}\mcite[def. IV.3.9]{Aluffi2009Algebra}
  We say that a group is \term[ru=разрешимая (группа) (\cite[48]{Кострикин2001АлгебраЧасть3})]{solvable} if its \hyperref[def:derived_group_series]{derived series} \hyperref[def:stabilizing_sequence]{stabilizes} at the trivial subgroup.
\end{definition}

\begin{proposition}\label{thm:abelian_group_is_solvable}
  Every \hyperref[def:abelian_group]{abelian group} is \hyperref[def:solvable_group]{solvable}.
\end{proposition}
\begin{proof}
  The commutator subgroup of an abelian group is trivial as shown in \cref{thm:abelian_group_iff_trivial_commutator}.
\end{proof}

\begin{proposition}\label{thm:simple_solvable_group_is_abelian}
  If a \hyperref[def:simple_group]{simple} group is \hyperref[def:solvable_group]{solvable}, it must be \hyperref[def:abelian_group]{abelian}.
\end{proposition}
\begin{proof}
  If \( G \) is simple, its derived group \( G' \) is either \( \set{ e } \) or \( G \).
  \begin{itemize}
    \item If \( G' = \set{ e } \), \cref{thm:abelian_group_iff_trivial_commutator} implies that \( G \) is abelian.
    \item If \( G' = G \), we can show by induction on \( k \) that \( G^{(k)} = G \). Unless \( G \) itself is the trivial group, in this case is not solvable because its derived series is constant.
  \end{itemize}
\end{proof}

\begin{example}\label{ex:def:solvable_group}
  We list examples of \hyperref[def:solvable_group]{solvable groups}:
  \begin{thmenum}
    \thmitem{ex:def:solvable_group/dihedral} We have shown in \cref{ex:def:group_commutator/dihedral} that the commutator of the dihedral group \( D_n \) is \( \braket{ a^2 } \).

    This subgroup is cyclic, hence abelian, so, by \cref{thm:abelian_group_is_solvable}, its own commutator is the trivial subgroup.

    We obtain the derived series
    \begin{equation*}
      \set{ e } \subsetneq \braket{ a^2 } \subsetneq D_n.
    \end{equation*}

    Therefore, \( D_n \) is solvable.

    \thmitem{ex:def:solvable_group/s3} We have shown in \cref{ex:def:group_commutator/s3} that the alternating group \( A_3 \) is the commutator of the symmetric group \( S_3 \).

    As in the previous example, it is abelian, so we have the derived series
    \begin{equation*}
      \set{ e } \subsetneq A_3 \subsetneq S_3.
    \end{equation*}

    Therefore, \( S_3 \) is solvable.

    \thmitem{ex:def:solvable_group/alternating} The alternating group \( A_n \) is simple when \( n \neq 4 \). This is shown in \cite[\S 9]{Курош1967ТеорияГрупп}.

    Additionally, \cref{thm:an_commutativity} implies that \( A_n \) is not abelian when \( n > 3 \).

    Therefore, \( A_5 \) is a simple non-abelian group, which is not solvable as as consequence of \cref{thm:simple_solvable_group_is_abelian}.
  \end{thmenum}
\end{example}
